\section*{第 \ref{ch:b3:lab-techniques-3} 章——实验技术 III：科研实践}

\begin{solution}{exo:b3:lab-techniques-3:1}
$(0.50/1013)^3 = \num{1.2e-10}$。泄漏、玻璃表面\omterm{def:b3:surfaces-catalysis:adsorption}{吸附}的气体及真空脂和隔垫的放气，
所遗留的量远大于此；干燥、加热过的玻璃器皿和良好的接口，
比多抽几次更重要。
\end{solution}

\begin{solution}{exo:b3:lab-techniques-3:2}
$\qty{20.0}{mmol}/\qty{1.48}{mol/L} = \qty{13.5}{mL}$.
\end{solution}

\begin{solution}{exo:b3:lab-techniques-3:3}
$120 + 10 \times 1.00782503 + 55.93493633 - 0.00054858 = \qty{186.01264}{u}$。误差：$10^6 \times (186.0129 - 186.01264)/186.01264 = +1.4$ ppm，
远在 5 ppm 之内。
\end{solution}

\begin{solution}{exo:b3:lab-techniques-3:4}
\omterm{def:b3:lab-techniques-3:literature}{一次文献}：期刊论文、专利、学位论文。\omterm{def:b3:lab-techniques-3:literature}{二次文献}：综述、
常数手册。三次文献：教科书。
\end{solution}

\begin{solution}{exo:b3:lab-techniques-3:5}
$M = \qty{194.0}{g/mol}$：C $96.0/194.0 = 49.48$\,\%，H $10.0/194.0 = 5.15$\,\%，N $56.0/194.0 = 28.87$\,\%。偏差分别为 0.17、0.07 和 0.16
个百分点，均在 0.4 之内，分析支持该化学式。
\end{solution}

\begin{solution}{exo:b3:lab-techniques-3:6}
$\qty{170.0}{mg}/\qty{212.0}{g/mol} = \qty{0.802}{mmol}$；$c = \qty{0.802}{mmol}/\qty{0.54}{mL} = \qty{1.5}{mol/L}$（两位有效数字，由所读
体积决定）。
\end{solution}

\begin{solution}{exo:b3:lab-techniques-3:7}
平均值分别为 73.0 和 77.0\,\%；两组的标准差均为 $\sqrt{10/4} = 1.58$；$s_p = 1.58$。$t = 4.0/(1.58\sqrt{2/5}) = 4.0$，大于
2.31，故该改进在 95\,\% 的水平上显著。
\end{solution}

\begin{solution}{exo:b3:lab-techniques-3:8}
$F = (0.80/0.30)^2 = 7.1 > 5.05$：第一种方法的精密度显著更低。
\end{solution}

\begin{solution}{exo:b3:lab-techniques-3:9}
购买带抑制剂的小瓶包装，收到和开封时均标注日期，贮存
于密闭、避光、阴凉处；任何蒸馏或蒸发前须检测过氧化物，
开封后每隔固定时间复测；超过有效期的瓶子，或
检测呈阳性的，应经危险废物途径处置，不得浓缩。
\end{solution}

\begin{solution}{exo:b3:lab-techniques-3:10}
危害：释放出的氢气（着火、增压），NaH 与水的反应活性，
以及已知的 DMF 与氢化钠的放热分解，其
加热时可能失控。DMF 还具有生殖毒性。控制措施：替换溶剂
（THF 或 2-甲基四氢呋喃）或替换碱；若仍需使用，先小量试做，低温下
分批加入 NaH 并监测温度，在通风橱内进行，
置于防护屏后，经鼓泡器排气，备好冷却浴；有书面操作规程，
有第二人在场，个人防护装备列于最后。
\end{solution}

\begin{solution}{exo:b3:lab-techniques-3:11}
相对贡献分别为 0.50、0.50、$0.010/20.0 = 0.05$、$0.010/15.0 = 0.067$ 和 0.10\,\%；合并后为
$\sqrt{0.25 + 0.25 + 0.0025 + 0.0044 + 0.010} = 0.72$\,\%。两项积分贡献最大：信噪比更高、
仔细的基线与相位校正，比天平精度更重要。
\end{solution}

\begin{solution}{exo:b3:lab-techniques-3:12}
$Q = (89.0 - 82.0)/(89.0 - 80.6) = 0.83 > 0.71$：检验判定其为可疑值。舍弃之前，先查看\omterm{def:b3:lab-techniques-3:notebook}{实验记录本}
看有无原因（如滴定管读数错误、气泡、滴定剂批次不同）；若找到
原因，则舍弃并说明理由；若未找到，则保留该值，报告检验结果，或
重新测量。
\end{solution}

\begin{solution}{pb:b3:lab-techniques-3:1}
\textbf{1.} $\ce{FeCl2 + 2NaC5H5 -> Fe(C5H5)2 + 2NaCl}$.

\textbf{2.} 环戊二烯基钠会被氧气和水破坏，而氯化亚
铁在潮湿空气中被氧化为三价铁。

\textbf{3.} $\qty{5.00}{g}/\qty{126.8}{g/mol} = \qty{0.0394}{mol}$；$\times \qty{185.8}{g/mol} = \qty{7.33}{g}$；产率 $5.86/7.33 = 80$\,\%。

\textbf{4.} $(0.10/1000)^3 = 10^{-12}$（理论上）。

\textbf{5.} 通过加热二聚体使其发生逆狄尔斯–阿尔德反应，并
蒸馏出挥发性单体；室温下单体会再次通过狄尔斯–阿尔德反应二聚，
因此须冷藏并立即使用。

\textbf{6.} 每一份都会释放氢气和热量：缓慢加入可使气体
流速与温度保持可控。

\textbf{7.} \qty{186.0126}{u}；误差 $+1.4$ ppm。

\textbf{8.} C $120.0/185.8 = 64.58$\,\%，H $10.0/185.8 = 5.38$\,\%；实测值为 64.41 和 5.47，偏差 $-0.17$ 和 $+0.09$，均在 0.4 之内。

\textbf{9.} 十个氢因对称性而等价，且环快速转动。

\textbf{10.} 只有一个信号，十个碳全部等价。

\textbf{11.} 单晶 X 射线结构分析。

\textbf{12.} 否：作为 18 电子配合物（\cref{ch:b3:organometallic-bonding}），它在空气中稳定，不同于其
前驱体。

\textbf{13.} 其信号与二茂铁的信号互不重叠，它是纯净、不挥发、
不吸湿的固体，可精确称量，且对
样品呈惰性。

\textbf{14.} 两次扫描之间每个信号充分弛豫（延迟数个 $T_1$），
激发均匀、信噪比足够，并仔细进行相位和基线
校正。

\textbf{15.} 二茂铁 $\ce{C10H10Fe}$ \qty{185.8}{g/mol}；1,3,5-三甲氧基苯 $\ce{C9H12O3}$ \qty{168.0}{g/mol}。

\textbf{16.} $P_x = 3.942 \times (3/10) \times (185.8/168.0) \times (15.0/20.0) \times 0.999 = 0.980$：98.0\,\%。

\textbf{17.} 相对不确定度 $\sqrt{0.50^2 + 0.50^2 + 0.05^2 + 0.067^2 + 0.10^2} = 0.72$\,\%；绝对不确定度 $0.72\,\% \times 98.0 = 0.7$ 个百分点。

\textbf{18.} 有影响，但很弱：与之组成相近的 2\,\% 杂质，会使 C 和 H 的测定值变化，
变化量小于\omterm{def:b3:lab-techniques-3:elemental}{元素分析}可检出的限度；qNMR 是信息量更大的纯度测定方法。

\textbf{19.} 氢化钠（遇水产生氢气，可引起火灾）、放出的氢气、THF 和
环戊二烯（高度易燃；THF 为\omterm{def:b3:lab-techniques-3:peroxide}{过氧化物形成剂}）、二聚体的热裂解、氯化铁(II)（刺激性）以及
二茂铁（易燃固体，有害）。

\textbf{20.} 氢化钠：在惰性气体保护下以油分散体形式称量，分批加入搅拌并冷却的溶液中，
产生的气体经鼓泡器排出，周围不得有水。THF：经柱子干燥并去除过氧化物，
在通风橱中、氮气保护下使用。

\textbf{21.} 在惰性气体保护并冷却条件下，缓慢加入异丙醇，随后加入乙醇，再加入水，
每次都须等到不再放出气体后再进行下一步。

\textbf{22.} 可逆的单电子波，$\ce{Fe(C5H5)2+}/\ce{Fe(C5H5)2}$，在 \qty{25}{\celsius} 下两峰间距约为 59 mV
（见\cref{ch:b3:electrode-kinetics}）；该电对反应快、稳定，且几乎不受
溶剂影响，因此可方便地用作内标。

\textbf{23.} \omterm{def:b3:lab-techniques-3:notebook}{实验记录本}的内容包括：日期、反应式、用量和来源、\omterm{def:b3:lab-techniques-3:assessment}{风险评估}、
所做和所见的内容、产率以及\omterm{def:b3:lab-techniques-3:notebook}{原始数据}的文件名。\omterm{def:b3:lab-techniques-3:literature}{补充材料}包括：
完整的实验步骤和全部表征数据，
并附谱图副本。

\textbf{24.} 以便他人，以及作者多年以后，无论使用何种软件，都能重新打开并再次处理这些数据，
即数据应可查找、可获取、可互操作、可再利用。

\textbf{25.} 二茂铁样品的 qNMR 纯度为 $98.0 \pm 0.7$\,\%（标准不确定度）。
\end{solution}
